\bprob[title={DoCarmo, Riemannian Geometry, page 103 question 1}]

    Let $G$ be a Lie group with a bi-invariant metric.
    Let $X,Y,Z\in\g$ be left-invariant vector fields.
    \benum
        \item Show that $\nabla_XY=\frac12[X,Y]$.
        \item Conclude that $R(X,Y)Z=\frac14[[X,Y],Z]$.
        \item Show that if $X,Y$ are orthonormal, and $\sigma$ is the plane generated by $X,Y$, then the sectional curvature of $G$ at $\sigma$ is
        $$ K(\sigma) = \frac12\norm{[X,Y]}^2 . $$
    \eenum

\eprob

\benum
    \item Consider the map $\g^2\to\g$ by $(X,Y)\mapsto\nabla_XY$.
    This is $\bR$-bilinear, as we know this of connections, and further we know that $\nabla_XX=0$.
    This implies that it is antisymmetric; a bilinear map $\alpha\colon V^2\to W$ which is zero when its inputs are equal is antisymmetric:
    $$ 0 = \alpha(v+u,v+u) = \alpha(v,v) + \alpha(v,u) + \alpha(u,v) + \alpha(u,u) = \alpha(v,u) + \alpha(u,v) . $$
    Thus $\nabla_XY=-\nabla_YX$ and by symmetry of the connection,
    $$ [X,Y] = \nabla_XY - \nabla_YX = 2\nabla_XY $$
    as desired.

    \item By the previous point,
    $$ R(X,Y)Z = \nabla_Y\nabla_XZ - \nabla_X\nabla_YZ + \nabla_{[X,Y]}Z = \frac14[Y,[X,Z]] - \frac14[X,[Y,Z]] + \frac12[[X,Y],Z] . $$
    Now by Jacobi,
    $$ 0 = [X,[Y,Z]] + [Z,[X,Y]] + [Y,[Z,X]] = [X,[Y,Z]] + [Z,[X,Y]] - [Y,[X,Z]] $$
    which implies
    $$ [Y,[X,Z]] - [X,[Y,Z]] = [Z,[X,Y]] = -[[X,Y],Z] . $$
    Thus we have
    $$ R(X,Y)Z = \frac14[[X,Y],Z] $$
    as desired.

    \item If $X,Y$ are orthonormal then the sectional curvature of $\sigma$ is just $\Rm(X,Y,X,Y)$ which is
    $$ K(\sigma) = \iprod{R(X,Y)X,Y} = \frac14\iprod{[[X,Y],X],Y} = -\frac14\iprod{\nabla_X[X,Y],Y} . $$
    Now, because the connection is metric,
    $$ X\iprod{[X,Y],Y} = \iprod{\nabla_X[X,Y],Y} + \iprod{[X,Y],\nabla_XY} . $$
    And similarly
    $$ 0 = X\iprod{Y,Y} = 2\iprod{\nabla_XY,Y} , $$
    where the left-hand side is zero since $\iprod{Y,Y}=1$ is constant.
    Thus $0=2\iprod{\nabla_XY,Y}=\iprod{[X,Y],Y}$, and so we see
    $$ 0 = \iprod{\nabla_X[X,Y],Y} + \iprod{[X,Y],\nabla_XY} \implies \iprod{[X,[X,Y]],Y} = -\iprod{[X,Y],[X,Y]} , $$
    and finally
    $$ K(\sigma) = \frac14\iprod{[X,Y],[X,Y]} $$
    as desired.
\eenum

\bprob

    \benum
        \item Let $\nabla$ be a connection on a manifold $M$.
        For vector fields $X,Y,Z\in\X(M)$, let $\nabla^2Z$ be the $(2,1)$ tensor field obtained by taking the covariant derivative of the $(1,1)$ tensor field $\nabla Z$.
        We write $\nabla^2_{X,Y}Z$ for $\nabla^2Z(X,Y)$.
        Show that
        $$ R(X,Y)Z = \nabla^2_{Y,X}Z - \nabla^2_{X,Y}Z . $$

        \item (DoCarmo, Riemmanian Geometry, page 106 problem 7): prove the second {\emphcolor Binachi identity}:
        $$ \nabla \Rm(X,Y,Z,W,T) + \nabla \Rm(X,Y,W,T,Z) + \nabla \Rm(X,Y,T,Z,W) = 0 . $$
    \eenum

\eprob

\benum
    \item By the formula for the covariant derivative of a tensor, we see that for a covector field $\omega$,
    $$ \multline{
        \nabla^2_{Y,X}Z(\omega) = \nabla_Y\nabla Z(\omega,X) = Y(\nabla Z(\omega,X)) - \nabla Z(\nabla_Y\omega,X) - \nabla Z(\omega,\nabla_YX)\cr
        = Y(\omega(\nabla_XZ)) - \nabla_Y\omega(\nabla_XZ) - \omega(\nabla_{\nabla_YX}Z) .
    } $$
    Now we know that in general,
    $$ Y(\omega X) = \nabla_Y(\omega X) = \nabla_Y\omega(X) + \omega(\nabla_YX) , $$
    and thus
    $$ Y(\omega(\nabla_XZ)) = \nabla_Y\omega(\nabla_XZ) + \omega(\nabla_Y\nabla_XZ) . $$
    So
    $$ \nabla^2_{Y,X}Z(\omega) = \omega(\nabla_Y\nabla_XZ) - \omega(\nabla_{\nabla_YX}Z) . $$
    Since this is true for all $\omega$, we see that
    $$ \nabla^2_{Y,X}Z = \nabla_Y\nabla_XZ - \nabla_{\nabla_YX}Z . $$
    Thus
    $$ \eqalign{
        \nabla^2_{Y,X}Z - \nabla^2_{X,Y}Z &= \nabla_Y\nabla_XZ - \nabla_{\nabla_YX}Z - \nabla_X\nabla_YZ + \nabla_{\nabla_XY}Z\cr
            &= \nabla_Y\nabla_XZ - \nabla_X\nabla_YZ + \nabla_{\nabla_XY-\nabla_YX}Z\cr
            &= \nabla_Y\nabla_XZ - \nabla_X\nabla_YZ + \nabla_{[X,Y]}Z\cr
            &= R(X,Y)Z.
    } $$

    \item Everything here is tensorial so we can choose whatever local frame we want, show it for that, and extrapolate the result for all vector fields.
    Let $p\in M$ and choose a geodesic field $(E_i)$ at $p$.
    Then we have,
    $$ \multline{
        \nabla\Rm(E_i,E_j,E_k,E_\ell,E_h) = \nabla_{E_h}\Rm(E_i,E_j,E_k,E_\ell) =\cr
            E_h(\Rm(E_i,E_j,E_k,E_\ell)) - \Rm(\nabla_{E_h}E_i,E_j,E_k,E_\ell) - \cdots - \Rm(E_i,\dots,\nabla_{E_h}E_\ell) .
    } $$
    At $p$ we know $\nabla_{E_i}E_j|_p=0$ for all $i,j$, so by multilinearity the last terms above are all zero (at $p$).
    So we see
    $$ \multline{
        \nabla\Rm(E_i,E_j,E_k,E_\ell,E_h) = E_h\Rm(E_i,E_j,E_k,E_\ell) = E_h\iprod{\nabla^2_{E_j,E_i}E_k-\nabla^2_{E_i,E_j}E_k,E_\ell} =\cr
        \iprod{\nabla_{E_h}(\nabla^2_{E_j,E_i}E_k-\nabla^2_{E_i,E_j}E_k),E_\ell} ,
    } $$
    the last equality is due to the connection being metric and $\nabla_{E_h}E_\ell=0$ at $p$.
    Now,
    $$ \iprod{\nabla_{E_h}\nabla^2_{E_j,E_i}E_k,E_\ell} = \iprod{\nabla_{E_h}\nabla_{E_j}\nabla_{E_i}E_k - \nabla_{E_h}\nabla_{\nabla_{E_j}E_i}E_k,E_\ell} , $$
    and so by this and symmetry of the connection,
    $$ \nabla\Rm(E_i,E_j,E_k,E_\ell,E_h) = \iprod{\nabla_{E_h}\nabla_{E_j}\nabla_{E_i}E_k - \nabla_{E_h}\nabla_{E_i}\nabla_{E_j}E_k + \nabla_{E_h}\nabla_{[E_i,E_j]}E_k,E_\ell} . $$

    Notice that $\Rm(E_i,E_j,E_k,E_\ell)=\Rm(E_k,E_\ell,E_i,E_j)$ so we also can write
    $$ \nabla\Rm(E_i,E_j,E_k,E_\ell,E_h) = \iprod{\nabla_{E_h}\nabla_{E_\ell}\nabla_{E_k}E_i - \nabla_{E_h}\nabla_{E_k}\nabla_{E_\ell}E_i + \nabla_{E_h}\nabla_{[E_k,E_\ell]}E_i,E_j} . $$
    And so we have that the sum gives the inner product of the following with $E_j$ (writing $\nabla_i$ for $\nabla_{E_i}$ and $\nabla_{i,j}$ for $\nabla_{[E_i,E_j]}$):
    $$ \displaylines{
        {}\nabla_h\nabla_\ell\nabla_kE_i - \nabla_h\nabla_k\nabla_\ell E_i + \nabla_h\nabla_{k,\ell}E_i\cr
        \mathllap{{}+} \nabla_\ell\nabla_k\nabla_hE_i - \nabla_\ell\nabla_h\nabla_kE_i + \nabla_\ell\nabla_{h,k}E_i\cr
        \mathllap{{}+} \nabla_k\nabla_h\nabla_\ell E_i - \nabla_k\nabla_\ell\nabla_h E_i + \nabla_k\nabla_{\ell,h}E_i .
    } $$
    Now consider
    $$ R(E_\ell,E_h)\nabla_kE_i = \nabla_h\nabla_\ell\nabla_kE_i - \nabla_\ell\nabla_h\nabla_kE_i + \nabla_{\ell,h}\nabla_kE_i . $$
    These are all zero at $p$ since $\nabla_kE_i=0$, and the first two terms are the terms from the the above large sum.
    The final term is zero since $[E_\ell,E_h]=\nabla_\ell E_h-\nabla_hE_\ell=0$, so the first two terms above cancel out, and so the large sum is just the sum of the terms of the form
    $\nabla_h\nabla_{k,\ell}E_i$.
    Also,
    $$ \nabla_h\nabla_{k,\ell}E_i = \nabla_h\nabla_{k,\ell}E_i - \nabla_{k,\ell}\nabla_hE_i = -R(E_h,[E_k,E_\ell])E_i + \nabla_{[E_h,[E_k,E_\ell]]}E_i , $$
    and since $[E_k,E_\ell]=0$ at $p$ we see that
    $$ \nabla_h\nabla_{k,\ell}E_i = \nabla_{[E_h,[E_k,E_\ell]]}E_i . $$
    Thus the large sum is the sum of terms of this form, as $h,k,\ell$ cycle.

    Thus the large sum is equal to
    $$ \sum_{(h,k,\ell)}\nabla_{[E_h,[E_k,E_\ell]]}E_i = \nabla_{\sum[E_h,[E_k,E_\ell]]}E_i , $$
    where the sum is over the cycles of $(h,k,\ell)$.
    By the Jacobi identity for the Lie bracket, this is zero.
    And thus we get that the sum of the covariant derivatives of $\Rm$ is the inner product of zero with $E_j$, which is zero.
\eenum

\bprob

    Let $E\to M$ be a vector bundle with metric $h$ and connection $\bar\nabla$.
    Let $S\subseteq E$ be a subbundle with orthogonal complement $S^\perp\subseteq E$ so that $E=S\oplus S^\perp$.
    Let $P^\perp\colon E\to S^\perp$ denote the orthogonal projection along $S$.
    Let $\nabla,\nabla^\perp$ denote the induced connections on $S,S^\perp$ respectively.
    Finally let $\sff\colon TM\otimes S\to S^\perp$ denote the second fundamental form of $S$ and $\bar R\colon\Lambda^2TM\otimes E\to E$ denote the curvature of $\bar\nabla$.
    \benum
        \item Show that for $X,Y\in\X(M)$ and $\xi\in\Gamma(S)$,
        $$ P^\perp\bar R(X,Y)\xi = \nabla^\perp_Y(\sff(X)\xi) - \sff(X)\nabla_Y\xi - \nabla^\perp_X(\sff(Y)\xi) + \sff(Y)\nabla_X\xi + \sff([X,Y])\xi . $$
        \item Let $M\subseteq\bar M$ be a submanifold, $\bar g$ a Riemannian metric on $\bar M$ and $g$ the induced metric on $M$.
        Let $E=T\bar M|_M$ and $S=TM$ so that $S^\perp=\nu_M$.
        Let $\bar\nabla$ be the connection induced from the Levi-Civita connection of $(\bar M,\bar g)$ so that $\nabla$ is the Levi-Civita connection on $(M,g)$.
        \benum
            \item Define $\nabla\sff\colon\X(M)\times\X(M)\times\X(M)\to\Gamma(\nu_M)$ by
            $$ \nabla_Y\sff(X)Z = \nabla^\perp_Y(\sff(X)Z) - \sff(\nabla_YX)Z - \sff(X)\nabla_YZ,\qquad X,Y,Z\in\X(M) . $$
            Show that $\nabla\sff$ is tensorial, and thus defines a bundle morphism
            $$ \nabla\sff\colon TM\otimes TM\otimes TM\to\nu_M . $$
            \item Show that
            $$ P^\perp\bar R(X,Y)Z = \nabla_Y\sff(X)Z - \nabla_X\sff(Y)Z . $$
        \eenum
    \eenum

\eprob

\benum
    \item The proof here is essentially the same as that of Gauss's theorem's proof from class.
    By definition we have for $X\in\X(M)$ and $\xi\in\Gamma(S)$,
    $$ \bar\nabla_X\xi = \nabla_X\xi + \sff(X)\xi . $$
    Since $\nabla_X\xi\in\Gamma(S)$ and $\sff(X)\xi\in\Gamma(S^\perp)$ repeating this we see
    $$ \bar\nabla_Y\bar\nabla_X\xi = \nabla_Y\nabla_X\xi + \sff(Y)\nabla_X\xi + \nabla^\perp_Y\sff(X)\xi + \sff^\perp(Y)\sff(X)\xi . $$
    The first and last terms here lie in $\Gamma(S)$, while the middle two lie in $\Gamma(S^\perp)$.
    Thus projecting along $S$ gives
    $$ P^\perp\bar\nabla_Y\bar\nabla_X\xi = \sff(Y)\nabla_X\xi + \nabla^\perp_Y\sff(X)\xi . $$
    And for the same reason as before, by definition,
    $$ \bar\nabla_{[X,Y]}\xi = \nabla_{[X,Y]}\xi + \sff([X,Y])\xi , $$
    and here the first term is in $\Gamma(S)$ while the second is in $\Gamma(S^\perp)$ so that
    $$ P^\perp\bar\nabla_{[X,Y]}\xi = \sff([X,Y])\xi . $$

    Putting this all together,
    $$ \eqalign{
        P^\perp\bar R(X,Y)\xi &= P^\perp(\bar\nabla_Y\bar\nabla_X\xi - \bar\nabla_X\bar\nabla_Y\xi + \bar\nabla_{[X,Y]}\xi)\cr
            &= \sff(Y)\nabla_X\xi + \nabla^\perp_Y\sff(X)\xi - \sff(X)\nabla_Y\xi - \nabla^\perp_X\sff(Y)\xi + \sff([X,Y])\xi,
    } $$
    which is the desired expression.

    \item
    \benum
        \item By adittivity of $\sff,\nabla,\nabla^\perp$ in all parameters, $\nabla\sff$ is clearly additive.
        And since $\sff$ is tensorial and $\nabla,\nabla^\perp$ are tensorial in their lower parameters, $\nabla\sff$ is tensorial in $Y$.
        It remains to show that $\nabla\sff$ is multuplicative in $X,Z$:
        $$ \nabla_Y\sff(fX)Z = f\nabla_Y\sff(X)Z,\qquad \nabla_Y\sff(X)(fZ) = f\nabla_Y\sff(X)Z\qquad\hbox{for $f\in C^\infty(M)$} . $$

        For $X$, we see
        $$ \eqalign{
            \nabla_Y\sff(fX)Z &= \nabla^\perp_Y(\sff(fX)Z) - \sff(\nabla_YfX)Z - \sff(fX)\nabla_YZ\cr
                &= \nabla^\perp_Y(f\sff(X)Z) - \sff((Yf)X + f\nabla_YX)Z - f\sff(X)\nabla_YZ\cr
                &= (Yf)\sff(X)Z + f\nabla^\perp_Y(\sff(X)Z) - (Yf)\sff(X)Z - f\sff(\nabla_YX)Z - f\sff(X)\nabla_YZ\cr
                &= f\bigl(\nabla^\perp_Y(\sff(X)Z) - f\sff(\nabla_YX)Z - f\sff(X)\nabla_YZ\bigr)\cr
                &= f\nabla_Y\sff(X)Z.
        } $$
        And similarly for $Z$, we see
        $$ \eqalign{
            \nabla_Y\sff(X)(fZ) &= \nabla^\perp_Y(\sff(X)(fZ)) - \sff(\nabla_YX)(fZ) - \sff(X)\nabla_Y(fZ)\cr
                &= \nabla^\perp_Y(f\sff(X)Z) - f\sff(\nabla_YX)Z - \sff(X)((Yf)Z + f\nabla_YZ)\cr
                &= (Yf)\sff(X)Z + f\nabla^\perp_Y(\sff(X)Z) - f\sff(\nabla_YX)Z - (Yf)\sff(X)Z - f\sff(X)\nabla_YZ\cr
                &= f\bigl(\nabla^\perp_Y(\sff(X)Z) - f\sff(\nabla_YX)Z - f\sff(X)\nabla_YZ\bigr)\cr
                &= f\nabla_Y\sff(X)Z.
        } $$
        So that $\nabla\sff$ is tensorial (linear) in all parameters.

        \item From our first point, it is sufficient to show that
        $$ \nabla_Y\sff(X)Z - \nabla_X\sff(Y)Z = \nabla^\perp_Y(\sff(X)Z) - \sff(X)\nabla_YZ - \nabla^\perp_X(\sff(Y)Z) + \sff(Y)\nabla_XZ + \sff([X,Y])Z . $$
        Direct computation yields
        $$ \nabla_Y\sff(X)Z - \nabla_X\sff(Y)Z = \nabla^\perp_Y(\sff(X)Z) - \sff(\nabla_YX)Z - \sff(X)\nabla_YZ - \nabla^\perp_X(\sff(Y)Z) + \sff(\nabla_XY)Z + \sff(Y)\nabla_XZ . $$
        Cancelling out similar terms we see that we need
        $$ \sff(\nabla_XY)Z - \sff(\nabla_YX)Z = \sff([X,Y])Z . $$
        And indeed, by tensorality of $\sff$ the left-hand side is equal to
        $$ \sff(\nabla_XY - \nabla_YX)Z = \sff([X,Y])Z , $$
        where the equality is due to symmetry of the Levi-Civita connection.
    \eenum
\eenum

\bprob[title={Do Carmo, Differential Geometry of Curves and Surfaces, q7 p169}]

    Consider the surface of revolution $(\phi(v)\cos u,\phi(v)\sin u,\psi(v))$ with constant Gaussian curvature $K$.
    Parameterize $\phi,\psi$ such that $(\phi')^2+(\psi')^2=1$.
    \benum
        \item Show that $\phi$ satisfies $\phi''+K\phi=0$ and $\psi$ is given by $\psi=\int\sqrt{1-(\phi')^2}$.
        \item Show that all surfaces of revolution with constant curvature $K=1$ which intersect the $xy$ plane orthogonally are given by
        $$ \phi(v) = C\cos v,\qquad \psi(v) = \int_0^v\sqrt{1-C^2\sin^2v} . $$
        for some constant $C$.
        \item Determine all surfaces of revolution with constant curvature $K=-1$, and show that they are all of one of the following forms
        $$ \phi(v) = e^v,\qquad \phi(v) = \cosh(v),\qquad \phi(v) = \sinh(v) . $$
        \item Show that the only surface of revolution with $K=0$ are the right cylinder, the right circular cone, and the pane.
    \eenum

\eprop

\benum
    \item We know that (Do Carmo page 162)
    $$ K = -\frac{\phi''}\phi , $$
    and thus $\phi''+K\phi=0$.
    And we know
    $$ (\psi')^2 = 1 - (\phi')^2 , $$
    we take the square root and integrate, giving us both equations.

    \item Here we need to solve the ODE $\phi''+\phi=0$.
    The characteristic polynomial of this ODE is $x^2+1=0$, which has roots $\pm i$.
    Taking the real and imaginary parts, we have solutions $\cos(v),\sin(v)$ and so the general solution is given by $A\cos(v)+B\sin(v)$.

    In order for the surface of revolution to intersect $xy$ orthogonally, all meridians must attain critical points when $v=0$.
    This is because geometrically, the meridian should be flat when intersecting the $xy$ plane, i.e.\ have a critical point.
    This just means $\phi'(0)=0$; so $0=-A\sin(0)+B\cos(0)=B$.

    So the surfaces of revolution are all of the form $\phi(v)=A\cos(v)$ and since $\phi'(v)=-A\sin(v)$,
    $$ \phi(v) = A\cos(v),\qquad \psi(v) = \int_0^v\sqrt{1-A^2\sin^2(v)} . $$

    \item We need to solve $\phi''-\phi=0$, which has characteristic polynomial $x^2-1$.
    The roots are $x=\pm1$, which means $\phi$ has general form $\phi(v)=Ae^v+Be^{-v}$.
    If $B=0$ we are finished, if $A=0$ then we can reparameterize $v$ as $-v$ to get $\phi(v)=Be^v$.
    Otherwise we claim that a shift of $v$ will give us $\phi(v)=C\cosh(v)$ or $C\sinh(v)$.
    So reparameterize by $v+c$, then
    $$ \phi(v)=Ae^{v+c}+Be^{-v-c} = Ae^c\cdot e^v + Be^{-c}\cdot e^{-v} . $$
    So we want to find a $c$ such that
    $$ \abs{Ae^c} = \abs{Be^{-c}} \iff e^{2c} = \abs{\frac BA} \iff c = \frac12\log\abs{\frac BA} . $$
    In such a case, letting $2C=Ae^c$ we see that
    $$ \phi(v) = Ce^v \pm Ce^{-v} , $$
    which is $C\cosh(v)$ or $C\sinh(v)$ (depending on sign), as desired.

    The formula for $\psi$ is due to $\cosh'(v)=\sinh(v)$ and $\sinh'(v)=-\cosh(v)$ and $(e^v)'=e^v$.

    \item We need to solve $\phi''=0$, which has solution $\phi(v)=Av+B$.
    Then we have $\psi(v)=Cv$ for $C=\sqrt{1-A^2}$.

    When $A=0$ this gives us
    $$ (B\cos u,B\sin u,v) $$
    which is the cylinder.
    When $A=1$ we have
    $$ ((v+B)\cos u,(v+B)\sin u,0) , $$
    which is the plane; we can shift $v$ to $v+B$ so this is just polar coordinates of the plane.
    In general for nonzero $A$ we can shift $v$ so that $B=0$.

    So for all other $A$, we have
    $$ (Av\cos u,Av\sin u,Cv) , $$
    and since $A\neq1$, $C\neq0$.
    This is the cone, as it is the stretching of $(v\cos u,v\sin u,v)$ which is the classic parameterization of the cone.
\eenum

\bprob

    \benum
        \item Show that any smooth manifold $M$ admits a Riemannian metric.
        \item Let $\nabla$ be the Levi-Civita connection on $M$.
        Define, for $\omega\in\Omega^k(M)$,
        $$ (\d\omega)(X_1,\dots,X_{k+1}) = \sum_{i=1}^{k+1}(-1)^{i+1}(\nabla_{X_i}\omega)(X_1,\dots,\hat X_i,\dots,X_{k+1}) . $$
        Show that this defines a linear operator
        $$ \d\colon\Omega^k(M) \to \Omega^{k+1}(M) . $$
        \item Show that in any coordinate frame, if
        $$ \omega = \sum_I\omega_I\d x^I \qquad\hbox{summing over increasing multi-indices} , $$
        then
        $$ \d\omega = \sum_I\d\omega_I\wedge\d x^I . $$
        \item Conclude that $\d$ does not depend on the underlying Riemannian metric, and that $\d$ coincides with the construction of the exterior derivative from class.
        \item Prove that
        $$ \multline{
            (\d\omega)(X_1,\dots,X_{k+1}) = \sum_{i=1}^{k+1}(-1)^{i+1}X_i(\omega(X_1,\dots,\hat X_i,\dots,X_{k+1}))\cr
            + \sum_{i<j}(-1)^{i+j}\omega([X_i,X_j],X_1,\dots,\hat X_i,\dots,\hat X_j,\dots X_{k+1}) .
        } $$
    \eenum

\eprob

\benum
    \item We can embed $M$ in $\bR^n$ for some $n$ (Whitney), and then pull back the Euclidean metric to $M$.
    Alternatively, cover $M$ by smooth charts $(U_i,\phi_i)$ where $U_i\subseteq M$ are open and $\phi_i\colon U_i\to\bR^n$ are diffeomorphisms onto open sets.
    For each $i$, pull back the Euclidean metric on $\bR^n$ to $U_i$ via $\phi_i$ to give $g_i=\phi_i^*\bar g$.
    This is a metric on $U_i$ as $\phi_i$ is a diffeomorphism.

    Let $(\psi_i)$ be a partition of unity subordinate to $(U_i)$, and define on all of $M$
    $$ g = \sum_i\psi_ig_i . $$
    Explicitly, we extend $g_i$ to be zero outside of $U_i$, and we define for $p\in M$, $g_p=\sum_i\psi_i(p)g_i|_p$.
    Around every $p\in M$ there is a neighborhood for which this sum is finite, and thus a finite sum of smooth tensor fields, so $g$ itself is a smooth tensor field.
    Furthermore it is a metric: $g_p$ is a linear combination of metrics $\psi_i(p)g_i|_p$ and is thus bilinear and symmetric.
    It is also nonnegative since each $g_i$ and $\psi_i$ are.
    And $g_p(v,v)=0$ iff $g_i|_p(v,v)=0$ for all $i$ where $p\in U_i$ and $\psi_i(p)>0$.
    There must be at least one such $i$, and since $g_i$ is a metric this means $v=0$, so $g$ is positive definite.

    \item We show note that $\d\omega$ is a tensor.
    We can write
    $$ \d\omega(X_1,\dots,X_{k+1}) = \sum_{i=1}^{k+1}(-1)^{i+1}\nabla\omega(X_1,\dots,\hat X_i,\dots,X_{k+1},X_i) , $$
    so
    $$ \d\omega = \sum\sgn\tau{}^\tau\nabla\omega, $$
    where the sum is over all permutations $\tau$ which take a specific index $i$ and move it to the end (meaning $i$ is mapped to $k+1$, and $j>i$ is shifted back one).
    Since $\nabla\omega$ is a tensor, permuting it is a tensor, and thus $\d\omega$ is a tensor.

    Now we need to show that $\d\omega$ is alternating.
    While $\nabla\omega$ itself is not alternating, $\nabla_X\omega$ is.
    This is because for any permutation $\sigma$,
    $$ \eqalign{
        \nabla_X\omega(Y_{\sigma1},\dots,Y_{\sigma k}) &= X(\omega(\bar Y_\sigma)) - \sum_i\omega(Y_{\sigma1},\dots,\nabla_XY_{\sigma i},\dots,Y_{\sigma k})\cr
            &= \sgn\sigma X(\omega(\bar Y)) - \sgn\sigma\sum_i\omega(Y_1,\dots,\nabla_XY_i,\dots,Y_k)\cr
            &= \sgn\sigma\nabla_X\omega(\bar Y).
    } $$
    Now, if $\bar X$ has a repeated value, say $X=X_j=X_{j+1}$,
    $$ \multline{
        \d\omega(\bar X) = \sum_i(-1)^{i+1}\nabla_{X_i}\omega(X_1,\dots,\hat X_i,\dots,X_{k+1}) =\cr
        (-1)^{j+1}\nabla_X\omega(\dots,\hat X_j,X,\dots) + (-1)^{j}\nabla_X\omega(\dots,X,\hat X_{j+1},\dots) ,
    } $$
    since all other terms in the sum have repeated inputs, and $\nabla_{X_i}\omega$ is alternating.
    Notice that there is an equality $(\dots,\hat X_j,X,\dots)=(\dots,X,\hat X_{j+1},\dots)$ and so this sum is zero.
    A tensor which is zero when it is input an adjacent repeated value is alternating.
    This is beacause then swapping adjacent values alters the sign: say $\alpha$ is multilinear and zero when adjacent values are repeated, then
    $$ \multline{
        0 = \alpha(\dots,X+Y,X+Y,\dots) = \alpha(\dots,X,X,\dots) + \alpha(\dots,X,Y,\dots) + \alpha(\dots,Y,X,\dots) + \alpha(\dots,Y,Y,\dots)\cr
        =\alpha(\dots,X,Y,\dots) + \alpha(\dots,Y,X,\dots) .
    } $$
    Since permutations of the form $(i,i+1)$ generate $S_n$, this means that such an $\alpha$ is then alternating.
    
    So $\d\omega$ is alternating, meaning is in $\Omega^{k+1}(M)$ as desired.
    Clearly by linearity of $\nabla\omega$ in $\omega$, $\d$ is linear.

    \item Let us denote by $D$ the linear map we just defined, and by $\d$ the exterior derivative.

    First we claim that $\nabla_X(\omega\wedge\eta)=\nabla_X(\omega)\wedge\eta+\omega\wedge\nabla_X\eta$ for all forms $\omega,\eta$ and $X\in\X(M)$.
    This is because $\nabla_X{}^\sigma F={}^\sigma\nabla_XF$ for all covariant tensors $F$ and permutations $\sigma$.
    This is immediate from
    $$ \nabla_XF(\bar Y) = X(F(\bar Y)) - \sum_iF(Y_1,\dots,\nabla_XY_i,\dots,Y_n) . $$
    Then $\nabla_X$ commutes with $\alt$ since it is just a linear combination of ${}^\sigma$s.
    So
    $$ \multline{
        \nabla_X(\omega\wedge\eta) = \nabla_X\alt(\omega\otimes\eta) = \alt\nabla_X(\omega\otimes\eta) = \alt((\nabla_X\omega)\otimes\eta) + \alt(\omega\otimes\nabla_X\eta)\cr
            = (\nabla_X\omega)\wedge\eta + \omega\wedge\nabla_X\eta .
    } $$

    Now we claim that $D(\d x^I)$ for any multi-index $I$ of length $k$.
    It is sufficient to show that for any multi-index $J$ of length $k+1$, $D(\d x^I)(\partial_J)=0$, where $\partial_J=(\partial_{j_1},\dots,\partial_{j_{k+1}})$.
    For a multi-index $I=(i)$ then we have
    $$ D(\d x^i)(\partial_j,\partial_\ell) = \nabla_{\partial_j}\d x^i(\partial_\ell) - \nabla_{\partial_\ell}\d x^i(\partial_j) . $$
    We know that for $1$-forms,
    $$ \nabla_{\partial_j}\d x^i(\partial_\ell) = \partial_\ell(\d x^i(\partial_j)) - \d x^i\nabla_{\partial_j}\partial_\ell = -\d x^i\nabla_{\partial_j}\partial_\ell , $$
    since $\d x^i\partial_j$ is constant.
    So this is equal to
    $$ D(\d x^i)(\partial_j,\partial_\ell) = \d x^i\bigl(\nabla_{\partial_\ell}\partial_j - \nabla_{\partial_j}\partial_\ell\bigr) = \d x^i[\partial_\ell,\partial_j] = 0 . $$
    So we have our base case.

    Now consider the multi-index $(I,\ell)$, where $I$ has length $k$.
    Writing $\hat J_i$ for the multi-index obtained by removing the index $i$ from $J$, we have
    $$ D(\d x ^I\wedge\d x^\ell)(\partial_J) = \sum_i(-1)^{i+1}\nabla_{\partial_{j_i}}\d x^I\wedge\d x^\ell(\partial_{\hat J_i}) + \sum_i(-1)^{i+1}\d x^I\wedge\nabla_{\partial_{j_i}}\d x^\ell(\partial_{\hat J_i}) . $$
    We claim that both sums here are zero.

    For the first sum, if $\ell\notin J$ then it is trivially zero.
    Otherwise, we can assume that it is the last index in $J$ (because we can otherwise permute $J$ and this will affect the sum only in sign), so let $J=(J',\ell)$.
    Then the first sum is
    $$ \sum_{i\leq k+1}(-1)^{i+1}\nabla_{\partial_{j_i}}\d x^I(\partial_{\hat J'_i}) = D(\d x^I)(\partial_{J'}) = 0 $$
    inductively.

    For the second sum, if $I$ is not a sub-multi-index of $J$ then this is trivially zero again.
    Otherwise we can assume $I$ occurs as the first $k$ indices of $J$, $J=(I,\partial_j,\partial_m)$ (otherwise we permute).
    Then in the sum, we can take $i$ to be greater than $k$ as otherwise we remove an index from $I$ in $J$ and the resulting wedge will be zero.
    So the sum is then equal to
    $$ \pm\bigl(\nabla_{\partial_j}\d x^\ell(\partial_m) - \nabla_{\partial_m}\d x^\ell(\partial_j)\bigr) = \pm D(\d x^\ell)(\partial_j,\partial_m) = 0 $$
    from our base case.
    So both sums are zero and thus inductively we have $D(\d x^I)=0$.

    By linearity of $D$ it is sufficient to show that $D(u\d x^I)=\d u\wedge\d x^I$ for $u\in C^\infty(M)$.
    Indeed, from what we just showed:
    $$ \eqalign{
        D(u\d x^I)(\partial_J) &= \sum_i(-1)^{i+1}\nabla_{\partial_{j_i}}u\d x^I(\partial_{\hat J_i})\cr
            &= \sum_i(-1)^{i+1}(\partial_{j_i}u)\d x^I(\partial_{\hat J_i}) + u\sum_i(-1)^{i+1}\partial_{j_i}\d x^I(\partial_{\hat J_i})\cr
            &= \sum_i(-1)^{i+1}\pdvof u{x_{j_i}}\delta^I_{\hat J_i} + uD(\d x^I)(\partial_J)\cr
            &= \sum_i(-1)^{i+1}\pdvof u{x_{j_i}}\delta^I_{\hat J_i}.
    } $$
    On the other hand,
    $$ \d u\wedge\d x^I(\partial_J) = \sum_i\pdvof u{x_i}\d x^i\wedge\d x^I(\partial_J) , $$
    and we can sum over $i$ occuring in $J$ as otherwise the term will be zero,
    $$ = \sum_i\pdvof u{x_{j_i}}\d x^{j_i}\wedge\d x^I(\partial_J) = \sum_i\pdvof u{x_{j_i}}\delta^{j_iI}_J . $$
    Since $j_i$ occurs at the $i$th index of $J$, a permutation of $J$ bring it to $j_iI$ must map $i$ to $1$ and gives a permutation of $\hat J_i$ mapping to $I$.
    The sign of this permutation is $(-1)^{i+1}$ times the sign of the sub-permutation mapping $\hat J_i$ to $I$, as we must offset all the indices before $i$ by $1$.
    So this is equal to
    $$ = \sum_i(-1)^{i+1}\pdvof u{x_{j_i}}\delta^I_{\hat J_i} = D(u\d x^I) , $$
    as desired.

    It is sufficient to show this equality for $\partial_J$ since they are both alternating and $\partial_J$ forms local frames.

    \item The formula
    $$ \d\omega = \sum_I\d\omega_I\wedge\d x^I $$
    is independent on choice of metric and connection, and it coincides with the construction from class.

    \item We know that
    $$ \nabla_Y\omega(\bar X) = Y(\omega(\bar X)) - \sum_i\omega(X_1,\dots,\nabla_YX_i,\dots,X_k) . $$
    Thus
    $$ \eqalign{
        \d\omega(X_1,\dots,X_{k+1}) &= \sum_i(-1)^{i+1}\nabla_{X_i}\omega(X_1,\dots,\hat X_i,\dots,X_{k+1})\cr
            &= \sum_i(-1)^{i+1}X_i(\omega(X_1,\dots,\hat X_i,\dots,X_{k+1})) - \sum_i(-1)^{i+1}\sum_{j\neq i}\omega(\dots,\hat X_i,\dots,\nabla_{X_i}X_j,\dots) .
    } $$
    In the final sum, $\hat X_i$ may come after $\nabla_{X_i}X_j$.
    This final sum can thus be written as
    $$ \sum_{i<j}(-1)^{j+1}\omega(\dots,\nabla_{X_j}X_i,\dots,\hat X_j,\dots) + \sum_{i<j}(-1)^{i+1}\omega(\dots,\hat X_i,\dots,\nabla_{X_i}X_j,\dots) . $$
    Permuting the inputs so that the covariant derivative appears first has a sign of $(-1)^{i+1}$ in the first sum and $(-1)^j$ in the second ($i-1$ and $j$ transpositions respectively).
    So this is equal to
    $$ = \sum_{i<j}(-1)^{i+j}\bigl(\omega(\nabla_{X_j}X_i,\dots) - \omega(\nabla_{X_i}X_j,\dots)\bigr) , $$
    and by symmetry of the connection, this is equal to
    $$ = \sum_{i<j}(-1)^{i+j}\omega([X_j,X_i],\dots) = -\sum_{i<j}(-1)^{i+j}\omega([X_i,X_j],\dots) . $$
    We subtract this from the first sum, so we see that
    $$ \d\omega(X_1,\dots,X_{k+1}) = \sum_i(-1)^{i+1}X_i(\omega(\dots\hat X_i\dots)) + \sum_{i<j}(-1)^{i+j}\omega([X_i,X_j],\dots) $$
    as desired.
\eenum



